\section{The nearby Euler characteristic along the boundary}
\label{sec:nearby-euler}

We now compare a nearby open fiber of the log model with the open
surfaces in its special fiber.  The essential local assertion is that
boundary intersections contribute zero, including the points where the
threefold is singular.  We first prove that assertion independently of
the global construction.

For a complex algebraic variety, write $\chi_c$ for the Euler
characteristic with compact supports.  On a complex algebraic variety it
equals the ordinary Euler characteristic.  If $f:Z\to\mathbb C$ is
regular, the \emph{nearby Euler function} $\psi_f$ is the unshifted
nearby-cycle operation on constructible functions.  In particular,
$\psi_f(\mathbf1_V)(z)$ is the Euler characteristic of the local Milnor
fiber of $f|_V$ when $V$ is closed.  Extend this definition by additivity
to constructible functions.  Equivalently, take stalk Euler
characteristics of the usual, unshifted sheaf nearby cycles.  This
operation depends only on the restriction away from $f^{-1}(0)$ and
commutes with proper pushforward.  For a finite map $\rho$, pushforward
is the sum over the finite set $\rho^{-1}(z)$; for a general proper map it
is integration with respect to $\chi_c$.  These conventions and the
proper compatibility are recalled in
\cite[Definition~3.2]{SchuermannNearby} and
\cite[Proposition~4.19(1)]{MaximSchuermann}.

\begin{lemma}[Boundary covers in dimension three]
\label{lem:boundary-cover}
Let $Z$ be a normal complex algebraic threefold and let
$B=\sum_{j=1}^n B_j$ be a reduced divisor with distinct prime components.
Assume that $(Z,B)$ is dlt and that every $B_j$ is $\mathbb Q$-Cartier.
If $z$ lies on at least two of the $B_j$, then, after replacing $Z$ by a
neighborhood $U$ of $z$, there is a finite surjective map
$\rho:\widetilde U\to U$ from a normal variety such that:
\begin{enumerate}
\item $\rho$ is \'etale in codimension one and over
      $U\setminus B$, and every total pullback $\rho^*B_j$ is a reduced
      Cartier divisor;
\item $(\widetilde U,\rho^*B)$ is dlt and is simple normal crossing at
      every point of $\rho^{-1}(z)$.
\end{enumerate}
\end{lemma}

\begin{proof}
Shrink $U$ so that the only boundary components present are those
through $z$, and choose regular functions $g_j$ with
$\operatorname{div}(g_j)=m_jB_j$, where $m_j$ is a positive integer.
Normalize $U$ in a field extension obtained by adjoining roots
$a_j^{m_j}=g_j$ for all $j$.  Normalization is finite here, and taking
one field extension gives a finite surjective map of some degree $d$.
This is the usual local boundary-Cartier cover; see also
\cite[Proposition~11.25]{Kollar2023Families}.

At a codimension-one point, write each $g_j$ as a unit times the
$m_j$-th power of a local equation, or as a unit if $B_j$ is absent.
Adjoining its root is therefore unramified there.  The same argument
shows that the cover is \'etale wherever all the $B_j$ are Cartier,
in particular over $U\setminus B$ and over the simple normal crossing
locus of $(U,B)$.  Consequently
$\operatorname{div}(a_j)=\rho^*B_j$ has coefficient one on every prime
in its support.  Thus the total pullback is reduced and Cartier.

Put $\widetilde B=\rho^*B$.  The codimension-one ramification formula
gives
\[
 K_{\widetilde U}+\widetilde B
   =\rho^*(K_U+B)
\]
with compatible canonical divisors.  More generally, if a divisorial
valuation $F$ upstairs restricts to $e$ times a normalized divisorial
valuation $E$ downstairs, then
\begin{equation}
\label{eq:finite-log-discrepancy}
 A(F,\widetilde U,\widetilde B)=eA(E,U,B),
\end{equation}
where $A$ denotes log discrepancy.  To see the formula, choose a
birational model on which $E$ is a divisor and normalize it in the
upper function field.  At the generic point of $F$, the ramification
divisor has coefficient $e-1$.  Hence the discrepancy coefficients
satisfy $a_F=e a_E+(e-1)$, giving
\eqref{eq:finite-log-discrepancy}.  Restriction of a divisorial
valuation through a finite extension is divisorial, since its value
group has finite index and its residue field has the same
transcendence degree.  This is the finite-cover discrepancy formula
of \cite[Proposition~5.20]{KM1998}; see also
\cite[11.22--11.23]{Kollar2023Families}.

It follows that the upper pair is log canonical and that every log
canonical center upstairs maps onto a log canonical center downstairs.
Each lower center meets the simple normal crossing locus because the
lower pair is dlt.  The cover is \'etale there, so each upper center
also meets the simple normal crossing locus.  This is the dlt
criterion.  We use the standard facts that the log canonical centers
of a dlt pair are the normal strata of its coefficient-one boundary,
with the expected codimensions, and that intersections of log
canonical centers are unions of such centers; see
\cite[Proposition~9.2 in the author draft]{FujinoDlt} and
\cite[Definition~11.5 and~11.10]{Kollar2023Families}.

Fix $y\in\rho^{-1}(z)$.  There is at least one upper boundary prime
through $y$ over each lower boundary prime through $z$.  Indeed each
root $a_j$ vanishes at $y$, so its Cartier divisor contains $y$.
Primes over distinct lower primes are distinct, since a finite map
preserves their dimensions.  Thus at least two upper primes meet at
$y$.  If at least three meet, the dlt stratum theorem makes $y$ a
zero-dimensional log canonical center.  The dlt criterion then says
that the upper pair is simple normal crossing at $y$.

Suppose instead that precisely two upper boundary primes $T_1,T_2$
meet at $y$.  They lie over distinct lower primes; hence each agrees
near $y$ with its entire reduced Cartier pullback and is itself
Cartier there.  The surface $T_1$ is normal by dlt adjunction, and a
normal surface is Cohen--Macaulay.  Therefore the Cartier curve
$C=T_2|_{T_1}$ has no embedded points.  It is generically reduced,
because the pair is simple normal crossing at the generic points of
its curve strata, so it is reduced.  Each irreducible component of
$C$ is a normal curve, hence smooth.  Two such components cannot meet
at $y$: their intersection would make $y$ a log canonical center,
forcing simple normal crossings there, whereas a simple normal
crossing threefold pair with only two boundary branches has no
zero-dimensional log canonical center.  Consequently $C$ is a smooth
curve near $y$.

Let $R=\mathcal O_{\widetilde U,y}$ and let $a,b$ define $T_1,T_2$.
Since $R/(a,b)$ is a regular local ring of dimension one,
\[
 \operatorname{edim}R
 \le \operatorname{edim}(R/(a,b))+2=3=\dim R.
\]
Thus $R$ is regular.  Moreover $a,b$ have independent images in its
cotangent space, so they are part of a regular system of parameters.
The two boundary components are transverse.  This proves the claimed
simple normal crossing assertion in every case.
\end{proof}

\begin{lemma}[Vanishing at boundary intersections]
\label{lem:boundary-nearby-zero}
Under the hypotheses of Lemma~\ref{lem:boundary-cover}, suppose that a
regular function $f$ has reduced zero divisor
$V=\sum_{j\in I}B_j$, where $I$ is a nonempty subset of
$\{1,\ldots,n\}$.  Put $H=\sum_{j\notin I}B_j$.
At every point $z\in V$ lying on at least two components of $B$,
\[
 \psi_f(\mathbf1_{Z\setminus H})(z)=0.
\]
\end{lemma}

\begin{proof}
Use the finite cover of Lemma~\ref{lem:boundary-cover}, and let $d$
denote its degree.  At each point
$y$ over $z$, choose analytic simple normal crossing coordinates.  As
the cover is unramified in codimension one, the vertical
multiplicities remain one.  After absorbing a unit into one
coordinate, the pulled-back function and horizontal divisor are
\[
 f\circ\rho=z_1\cdots z_r,
 \qquad
 \rho^{-1}H=\{z_{r+1}\cdots z_{r+h}=0\},
 \qquad r\ge1,\quad r+h\ge2.
\]
When $h=0$, the horizontal divisor is empty.  For $r\ge2$, the local
Milnor fiber of the coordinate product has the homotopy type of
$(S^1)^{r-1}$ and therefore Euler characteristic zero.  The same is
true after restricting to any intersection of horizontal coordinate
hyperplanes.  Inclusion--exclusion therefore gives zero for the
complement of the horizontal boundary.  For $r=1$, each such
restricted Milnor fiber is a polydisc, with Euler characteristic
one.  Here $h\ge1$, so inclusion--exclusion gives
$\sum_{J\subseteq\{1,\ldots,h\}}(-1)^{|J|}=(1-1)^h=0$.
These computations may be made on small adapted polydiscs, or on
Milnor balls, with the same local nearby Euler values.

Let $\widetilde\alpha=\mathbf1_{\widetilde U\setminus\rho^{-1}H}$.
Over $U\setminus V$, its finite pushforward is
$d\mathbf1_{U\setminus H}$: away from $H$ the cover is \'etale of
degree $d$, and on $H$ all the terms are zero.  Since nearby cycles
only use the restriction off $V$, proper compatibility gives
\[
 d\,\psi_f(\mathbf1_{U\setminus H})(z)
  =\sum_{y\in\rho^{-1}(z)}
    \psi_{f\circ\rho}(\widetilde\alpha)(y)=0.
\]
This is an identity of constructible Euler functions; no triviality
of the covering local system is needed.  Since $d>0$, the result
follows.
\end{proof}

\begin{proposition}[Euler comparison for the log model]
\label{prop:nearby-euler-comparison}
Let $q:Y\to C$ be a projective flat morphism from a normal complex
threefold to a smooth pointed curve $(C,0)$.  Let $H$ be a reduced
divisor with horizontal prime components.  Assume that the Cartier
fiber $Y_0=\sum S$ is reduced, that $(Y,Y_0+H)$ is dlt, and that its
boundary components are $\mathbb Q$-Cartier.  Assume also that
$Y_t\setminus H$ is smooth for all sufficiently small $t\ne0$.
For each component $S$ of $Y_0$, let $B_S$ be the reduced union of its
intersections with the other components of $Y_0+H$, and put
$S^\circ=S\setminus B_S$.  If each local smoothing of an isolated
singularity in $S^\circ$ has Milnor fiber Euler characteristic at
least one, then
\begin{equation}
\label{eq:nearby-euler-comparison}
 \chi_c(Y_t\setminus H)
   =\sum_{S\subset Y_0}\chi_c(S^\circ)+\delta,
 \qquad \delta\ge0,
\end{equation}
for all sufficiently small $t\ne0$.
\end{proposition}

\begin{proof}
Write $\lambda=\psi_q(\mathbf1_{Y\setminus H})$ using a local
coordinate on $C$ at $0$.  Proper compatibility gives
\[
 \chi_c(Y_t\setminus H)=\int_{Y_0}\lambda\,d\chi_c.
\]
At every point lying on at least two components of $Y_0+H$,
Lemma~\ref{lem:boundary-nearby-zero} gives $\lambda=0$.
Every remaining point lies in exactly one $S^\circ$.  The surface
$S$ is normal by dlt adjunction, so its singular points form a finite
set.  At a smooth point of $S^\circ$, the local reduced Cartier fiber
is precisely $S$.  Its regular local ring is a quotient of the
ambient local ring by the base parameter, so the ambient ring is
regular and that parameter is part of a regular system of
parameters.  Thus $q$ is smooth there and $\lambda=1$.

At a singular point $p\in S^\circ$, a neighborhood misses $H$ and all
other special-fiber components.  Flatness and the hypothesis on the
nearby open fibers therefore give a smoothing of the normal isolated
surface singularity $(S,p)$.  If $M_p$ is its Milnor fiber, then
$\lambda(p)=\chi(M_p)\ge1$ by the assumed interior smoothing
statement.  Integrating the disjoint open surfaces and these finitely
many corrections yields \eqref{eq:nearby-euler-comparison}, with
\[
 \delta=\sum_{S\subset Y_0}
        \sum_{p\in\operatorname{Sing}(S)\cap S^\circ}
             \bigl(\chi(M_p)-1\bigr)\ge0.
\]
\end{proof}

In the projective log model $h:Y\to X$, the special-fiber components
are $S_0$ and the extra surfaces $S_a$.  The preceding proposition
applies once the interior smoothing inequality is supplied.  In
particular, the boundary calculation does not require the different
obtained by further adjunction to a curve to have integral
coefficients.
